En el fons d'una piscina de 2~m de profunditat hi ha una llanterna encesa, que emet llum en totes les direccions. La piscina està tapada amb una lona que queda just per sobre de l'aigua i en la lona s'observa un cercle de llum a causa dels raigs refractats. Això es deu a l'existència d'un angle límit, a partir del qual es dona el fenomen de reflexió total. Els rajos que arriben a la superfície amb un angle més gran que l'angle límit són plenament reflectits, i no emergeixen a la superfície. L'índex de refracció de l'aire és $n_0=1$ i de l'aigua és $n_a=1,33$.

\figura[0.35]{fig1.png}

\begin{apartat}{1,25}
A partir de la llei d'Snell, deduïu l'expressió de l'angle límit. Expresseu el resultat en funció de l'índex de refracció de l'aigua i de l'aire. Justifiqueu si podria donar-se el mateix fenomen si la situació fos la contrària (la llanterna estigués situada 2~m per sobre la superfície de la piscina i la lona dins de l'aigua just per sota la superfície).
\end{apartat}

\begin{apartat}{1,25}
Calcula el radi del cercle de llum, sabent que l'índex de refracció de l'aigua és $n=1,33$
\end{apartat}
